Welcome to the 2026 Widening NLP (WiNLP) Workshop!

\vspace{\baselineskip}

We are delighted to present this year’s proceedings and to welcome our community to another edition of Widening NLP. The workshop provides a space for researchers from underrepresented backgrounds to share their work, exchange ideas, and build connections across the Natural Language Processing (NLP) community.

\vspace{\baselineskip}

Widening NLP grew out of a conversation at ACL 2016 in Berlin, where a small group of researchers came together to discuss the underrepresentation of women and other minorities in NLP. That conversation led to the inaugural Workshop for Women and Underrepresented Minorities in NLP at ACL 2017. Since then, the workshop has continued to evolve, guided by the conviction that a stronger NLP community depends on making room for a wider range of voices, experiences, and perspectives.

\vspace{\baselineskip}

This commitment has shaped both the scope and the structure of the workshop. Following the inaugural event, we introduced two submission deadlines, including an early deadline to give authors additional time for visa applications. In 2019, we broadened our emphasis on diversity to encompass scientific and disciplinary backgrounds, educational pathways, training, and career stages. We also introduced a peer feedback system, allowing authors to receive constructive input from colleagues before formal review. Together, these efforts have sought to reduce barriers to participation and create a supportive environment for sharing research.

\vspace{\baselineskip}

The 2026 program continues this tradition. This year, WiNLP places particular emphasis on accessibility and disability inclusion. We also continue to welcome researchers from diverse scientific and disciplinary backgrounds and training pathways, as well as research on underrepresented languages. We received 67 submissions and accepted 16 archival and 9 non-archival papers. We are pleased to feature invited talks by XXX, XXX, and XXX. Their contributions, together with the work of our authors, offer opportunities to engage with diverse research perspectives and reflect on the future of NLP.

\vspace{\baselineskip}

We thank our authors, reviewers, invited speakers, panelists, and participants for the time, care, and expertise they have contributed to this workshop. We hope that the research and conversations shared here will inspire new ideas, strengthen collaborations, and support continued efforts toward a more inclusive and welcoming NLP community.

\vspace{\baselineskip}

-Chen, Emily, Meryem, Peerat, Richard, Santosh, Surendrabikram, Yinqiao, Organizing Co-Chairs